% !TeX root = ../main.tex
% !TeX encoding = UTF-8

\chapter*{Appendix A: The Governance Classification, Complete Framework}
\addcontentsline{toc}{chapter}{Appendix A: The Governance Classification, Complete Framework}
\markboth{Appendix A: The Governance Classification, Complete Framework}{Appendix A: The Governance Classification, Complete Framework}

The Governance Classification locates an organisation's AI governance position across five independent facets: Accountability, Exposure, Control, Regulation, and Maturity. Each facet is assessed at one of five levels. The result is a governance address, expressed as A{[}1-5{]}·E{[}1-5{]}·C{[}1-5{]}·R{[}1-5{]}·M{[}1-5{]}.

The framework is derived from S.R. Ranganathan's Colon Classification theory (1933), which established that complex subjects cannot be reduced to a single hierarchy and must be expressed as a combination of independent facets. The application of faceted classification to AI governance accountability is original to Otopoetic Limited.

The address is not a score. The five facets are independent, and no level on one facet compensates for a level on another. An organisation at A4·E1·C3·R2·M3 has named board-level accountability with no systematic AI inventory. A composite score would average that position into apparent adequacy. The address keeps the gap structurally visible.

\section{The Five Facets and Their Levels}

\textbf{A, Accountability.} The presence and depth of named individual ownership of AI systems and decisions at board level. The facet most directly engaged by the SM\&CR Duty of Responsibility and the Companies Act 2006 section 174 duty of care.

\begin{itemize}
\tightlist
\item
  A1. No documented ownership of AI systems or decisions at any level.
\item
  A2. Ownership assigned informally or at team level; no board visibility.
\item
  A3. Executive accountability exists; board ownership incomplete.
\item
  A4. Named board-level accountability for each material AI system, documented and current.
\item
  A5. The full Digital Alibi: all four elements present and active at the output-moment level for every material AI-assisted decision, with information provenance, a retrieval log, a human observability record, and a retrievability guarantee.
\end{itemize}

\textbf{E, Exposure.} The completeness and currency of the AI systems inventory, including third-party and vendor-embedded AI. The facet most directly engaged by DORA's ICT risk requirements and the EU AI Act's obligation to identify all high-risk systems in scope.

\begin{itemize}
\tightlist
\item
  E1. No systematic inventory; shadow AI likely present.
\item
  E2. Production systems mapped; third-party and shadow AI unmonitored.
\item
  E3. All AI systems identified; risk not formally assessed.
\item
  E4. Complete inventory with risk classification, refreshed on a defined cycle.
\item
  E5. Exposure profile updates continuously, with emerging risks proactively identified.
\end{itemize}

\textbf{C, Control.} The operational state and independence of governance structures at the moment of decision. The facet most directly engaged by ISO/IEC 42001 certification and the EU AI Act's quality management requirements.

\begin{itemize}
\tightlist
\item
  C1. Governance exists on paper only.
\item
  C2. Controls defined; operation inconsistent and untested.
\item
  C3. Governance framework operational; not independently tested.
\item
  C4. Controls independently verified; evidence retained for regulators.
\item
  C5. Governance evolves with the AI landscape; board assurance is a standing operational discipline.
\end{itemize}

\textbf{R, Regulation.} The completeness of the organisation's identification of, and compliance with, the AI governance obligations that apply to it. The facet most directly affected by the EU AI Act compliance clock.

\begin{itemize}
\tightlist
\item
  R1. Applicable obligations not identified or assessed.
\item
  R2. Obligations identified; gap analysis incomplete.
\item
  R3. Compliance programme active; not yet audit-ready.
\item
  R4. Obligations materially met, documented, and defensible to regulators.
\item
  R5. Ahead of current obligations; positioned for emerging requirements before they are mandated.
\end{itemize}

\textbf{M, Maturity.} The depth and direction of governance development as operational reality rather than documentation.

\begin{itemize}
\tightlist
\item
  M1. No structured board view of AI governance requirements.
\item
  M2. Awareness established; work not yet resourced.
\item
  M3. Governance work underway; not yet embedded organisationally.
\item
  M4. Governance embedded and tested under operational pressure.
\item
  M5. Governance capability is a competitive and reputational asset, demonstrable in due diligence and board assurance.
\end{itemize}

\section{Interpreting the Address}

Three positions carry particular significance. First, A2·E2·C2·R2·M2, the voluntary principles position, provides no legal protection. When an incident occurs, voluntary commitments are adduced as evidence that the organisation knew the standard it had not met. Second, A3·E3·C3·R3·M3 is the convergence point at which the major formal compliance standards engage; below it on any facet, the organisation has not met any formal standard. Third, A4 across all facets is the formally verified position that the EU AI Act's high-risk obligations, ISO/IEC 42001, and DORA collectively require, and A5 is the evidentiary standard beyond it: the address at which the full Digital Alibi is active. No existing regulatory instrument requires A5. A regulatory examination, litigation discovery, or shareholder challenge will.

\section{Self-Assessment Method}

Assess each facet independently, in writing, against the level descriptors, citing the evidence that supports the level claimed. Where the evidence for a level cannot be produced within the timeframe an examiner would allow, record the facet at the level the producible evidence supports. The gap between the self-assessed address and the confirmed address, established through independent examination of specific decisions, is itself a governance finding, and it is typically largest on the E facet, where shadow AI and vendor-embedded AI sit outside the formal inventory.
